\section{Implementation Details}
\label{app:implementation}

This appendix records the implementation specifics that the methods of
Section~\ref{sec:methods} depend on but that would interrupt the argument in the
main text: the experience-replay variants, and the NCL realisation together with
the hyperparameters fixed by the rot-MNIST tuning sweep.

\subsection{Datasets}
\label{app:impl_data}

rot-MNIST rotates each $28\times28$ image by the per-task angle and feeds the
flattened greyscale vector to the MLP. The CIFAR-10 domain shifts are generated
with the \texttt{imagecorruptions} package, which implements the standard
common-corruptions benchmark of \citet{hendrycks2019robustness}; the Gaussian
noise, shot noise, and contrast shifts are applied at severity $3$ on the
benchmark's $1$--$5$ scale, per sample at load time, followed by the standard
CIFAR-10 channel normalisation. The clean task (``none'') applies the
normalisation only.

\subsection{Models and Training}
\label{app:impl_models}

Tables~\ref{tab:mlp} and~\ref{tab:resnet} give the two backbones, and
Table~\ref{tab:training} the training protocol shared across both benchmarks.

\begin{table}[H]
    \centering
    \small
    \begin{tabular}{@{}l p{0.55\columnwidth}@{}}
        \hline
        Component        & Value                                    \\
        \hline
        Input            & $784$ (flattened $28\times28$ greyscale) \\
        Hidden layers    & $2 \times 400$ (fully connected, ReLU)   \\
        Output           & $10$-class logits                        \\
        Regularisation   & none                                     \\
        Total parameters & $478{,}410$                              \\
        \hline
    \end{tabular}
    \caption{MLP backbone used for the rot-MNIST sweeps.}
    \label{tab:mlp}
\end{table}

\begin{table}[H]
    \centering
    \small
    \begin{tabular}{@{}l p{0.55\columnwidth}@{}}
        \hline
        Component        & Value                                                                                               \\
        \hline
        Input            & $3\times32\times32$ RGB                                                                             \\
        Stem             & $3\times3$ conv, stride $1$, no max-pool (CIFAR adaptation)                                         \\
        Stages           & $4$ residual stages, BasicBlocks $[2,2,2,2]$, widths $[64,128,256,512]$                             \\
        Normalisation    & BatchNorm (default); GroupNorm ($\min(32, C)$ groups) variant crossed in Section~\ref{subsec:cifar} \\
        Head             & global average pool $+$ linear, $10$-class logits                                                   \\
        Total parameters & ${\approx}11.2$M                                                                                    \\
        \hline
    \end{tabular}
    \caption{CIFAR-adapted ResNet-18 backbone used for the CIFAR-10 block. The $3\times3$ stride-$1$ stem and removed max-pool keep the feature map at $32\times32$; all other details follow the ResNet-18.}
    \label{tab:resnet}
\end{table}

\begin{table}[H]
    \centering
    \small
    \begin{tabular}{@{}l p{0.27\columnwidth} p{0.27\columnwidth}@{}}
        \hline
        Hyperparameter  & rot-MNIST (MLP)         & CIFAR-10 (ResNet-18)   \\
        \hline
        Epochs per task & $1$ (online)            & $10$                   \\
        Batch size      & $256$                   & $256$                  \\
        Optimiser       & SGD                     & SGD                    \\
        Learning rate   & $0.1$                   & $0.1$                  \\
        Momentum        & $0.0$ or $0.9$          & $0.9$ ($0.0$ in cross) \\
        Weight decay    & $0.0$                   & $0.0$                  \\
        Replay buffer   & $1{,}000$ or $60{,}000$ & $5{,}000$              \\
        Seeds           & $5$                     & $5$ or $3$             \\
        \hline
    \end{tabular}
    \caption{Training protocol for both benchmarks.}
    \label{tab:training}
\end{table}

\subsection{Experience replay}
\label{app:impl_er}

The replay buffer uses reservoir sampling \citep{vitter1985random}, so the stored
set is a uniform sample of the past stream at any budget. Three modes appear in
the experiments:

\begin{itemize}
    \item \emph{Standard ER} (the G1 and all C-series conditions): each step takes
          one combined backward pass on $L_{\text{current}} + L_{\text{replay}}$,
          with half the mini-batch drawn from the current task and half sampled
          with replacement from the buffer.
    \item \emph{Balanced ER} (G3): two separate backward passes whose gradients are
          individually unit-normalised before being combined $50/50$, removing the
          magnitude asymmetry between $g_{\text{new}}$ and $g_{\text{replay}}$.
    \item \emph{Full-data replay} (G2, G4): the replay gradient is computed on
          \emph{every} stored sample exactly once per step (no sub-sampling). Paired
          with a buffer equal to the full past-task training set ($60{,}000$ for
          rot-MNIST $T_0$), this yields the deterministic empirical past-task
          gradient at the current parameters, removing estimator variance.
\end{itemize}

The $\lambda$-curriculum re-weights only the current-task term, minimising
$L_\lambda = L_{\text{replay}} + \lambda(t)\, L_{\text{current}}$; the adaptive
schedule reads $\|g_{\text{replay}}\|/\|g_{\text{new}}\|$ through an exponential
moving average ($\alpha = 0.05$) and applies the optional floor $\lambda_{\min}$.

\subsection{Asymmetric PER}
\label{app:impl_per}

The filtered push $\delta\,(F_{\text{rep}}+\delta I)^{-1} g_{\text{cur}}$ of
Equation~\eqref{eq:asym_per} is computed each step by warm-started conjugate
gradients over Fisher-vector products ($25$ iterations), so the metric is never
formed as a matrix; warm-starting from the previous step's solution keeps the
per-step solve cheap across the run. The solver control of
Section~\ref{subsec:per_asym} reruns the strongest filter ($\delta = 0.03$, full
buffer) at $60$ iterations to separate properties of the gate from artefacts of
an under-converged solve. On the full-buffer leg the replay gradient is exact
over all $60$k past samples while the Fisher stays estimated on a sampled batch,
since a full-buffer Fisher-vector product per CG iteration would dominate the
runtime.

\subsection{Natural Continual Learning}
\label{app:impl_ncl}

NCL follows \citet{kao2021naturalcontinuallearningsuccess} (their Eq.~8 and
Algorithm~1). The update preconditions the new-task gradient by the prior
covariance $\Lambda_{k-1}^{-1}$ and adds the Bayesian restoring term,
\begin{equation*}
    \theta \leftarrow \theta - \eta \left[ \Lambda_{k-1}^{-1} \nabla L_k(\theta) + (\theta - \mu_{k-1}) \right],
\end{equation*}
with $\Lambda_{k-1}$ summarised by a K-FAC factorisation $F \approx A \otimes G$
per linear layer and an identity prior $p_w = \alpha\, I$. The trust radius of the
paper's Eq.~7 is absorbed into the learning rate; we log a \texttt{tr\_scale}
diagnostic but do not clip the step.

\paragraph{Tuned hyperparameters.} The values used in every NCL condition are
those that won the rot-MNIST tuning sweep below (Table~\ref{tab:ncl_hparams}). The
paper-literal prior $\alpha = 1.0$ over-regularises on rot-MNIST and the damping
$\varepsilon$ is a numerical safeguard only: the identity prior, not the
damping, is what bounds $\Lambda^{-1}$. These values were tuned on the two-task rot-MNIST/MLP setting and reused unchanged on the CIFAR-10 ResNet-18, as we did not run a separate NCL sweep on the deeper backbone. NCL's CIFAR results should therefore be read as those of a method tuned on the smaller benchmark rather than re-optimised for CIFAR, a caveat we revisit in Section~\ref{subsec:disc_limits}.

\begin{table}[H]
    \centering
    \small
    \begin{tabular}{@{}l l l@{}}
        \hline
        Hyperparameter            & Symbol                           & Value             \\
        \hline
        Prior precision           & $\alpha$ (\texttt{prior\_init})  & $0.1$             \\
        Damping                   & $\varepsilon$ (\texttt{damping}) & $1\mathrm{e}{-3}$ \\
        Fisher samples (K-FAC)    & \texttt{fisher\_samples}         & $1{,}000$         \\
        Trust radius (diagnostic) & $r$ (\texttt{trust\_radius})     & $1.0$             \\
        Learning rate             & $\eta$                           & $0.1$             \\
        Momentum                  & $\mu$                            & $0.9$             \\
        \hline
    \end{tabular}
    \caption{NCL hyperparameters fixed for all conditions, from the rot-MNIST
        tuning sweep of Table~\ref{tab:ncl_sweep}.}
    \label{tab:ncl_hparams}
\end{table}

\paragraph{Tuning sweep.} The prior precision $\alpha$ was tuned on the two-task
rot-MNIST/MLP setting (Table~\ref{tab:ncl_sweep}), holding damping at
$1\mathrm{e}{-3}$ and \texttt{fisher\_samples} at $1{,}000$. Smaller $\alpha$
improves over the paper-literal $\alpha = 1.0$, but the gain saturates: $\alpha =
    0.1$ gives the best stability--accuracy trade-off, while $\alpha \leq 0.03$ becomes
unstable at $\eta = 0.1$ (the natural gradient amplifies the raw gradient by up to
$1/\alpha$ in any direction), and halving the learning rate does not recover it.

\begin{table}[H]
    \centering
    \small
    \begin{tabular}{@{}l l l l@{}}
        \hline
        Config                       & $\alpha$ & $\eta$ & Momentum \\
        \hline
        \textbf{alpha\_0.1} (chosen) & $0.1$    & $0.1$  & $0.9$    \\
        alpha\_0.03                  & $0.03$   & $0.1$  & $0.9$    \\
        alpha\_0.01                  & $0.01$   & $0.1$  & $0.9$    \\
        alpha\_0.003                 & $0.003$  & $0.1$  & $0.9$    \\
        alpha\_0.03\_lowlr           & $0.03$   & $0.05$ & $0.9$    \\
        \hline
    \end{tabular}
    \caption{NCL prior-precision tuning sweep on two-task rot-MNIST. Damping
        ($1\mathrm{e}{-3}$) and \texttt{fisher\_samples} ($1{,}000$) held fixed.}
    \label{tab:ncl_sweep}
\end{table}
